\chapter{Le VIH/SIDA : infection, prévention, sensibilisation}

\noindent\textit{Le VIH est un virus qui s'attaque directement au système immunitaire
étudié dans les deux chapitres précédents. Comprendre son fonctionnement, ses modes de
transmission réels et les moyens de s'en protéger est essentiel pour se prémunir soi-même
et lutter contre les idées fausses qui entourent cette maladie.}
\vspace{8pt}

\connecteBanner

\section{Le VIH, un virus qui cible l'immunité}

\begin{definition}[VIH]
Le \textbf{VIH} (Virus de l'Immunodéficience Humaine) est un virus qui infecte et détruit
progressivement les \textbf{lymphocytes T}, des cellules essentielles de l'immunité
acquise (voir chapitre précédent).
\end{definition}

\begin{propriete}[Conséquence]
En détruisant peu à peu les lymphocytes T, le VIH \textbf{affaiblit} progressivement le
système immunitaire de la personne infectée, la rendant de plus en plus vulnérable aux
infections.
\end{propriete}

\section{L'évolution de l'infection}

\begin{propriete}[Trois phases]
\begin{itemize}
\item \textbf{Contamination et séroconversion} : après la contamination, l'organisme
produit des anticorps anti-VIH, détectables par un test quelques semaines plus tard : la
personne devient \textbf{séropositive}.
\item \textbf{Phase asymptomatique} : le virus continue à détruire des lymphocytes T,
lentement, souvent \textbf{pendant plusieurs années}, sans symptôme apparent. La personne
séropositive se sent en bonne santé mais peut transmettre le virus.
\item \textbf{SIDA déclaré} : lorsque le nombre de lymphocytes T devient trop faible, le
système immunitaire s'effondre et des \textbf{maladies opportunistes} (infections,
cancers, que l'organisme ne parvient plus à combattre) peuvent se développer.
\end{itemize}
\end{propriete}

\begin{illus}
\begin{tikzpicture}[scale=0.8]
\draw[->,onbmuted] (0,0) -- (10.5,0) node[right,font=\small] {temps};
\draw[->,onbmuted] (0,0) -- (0,4.6) node[above,font=\small] {lymphocytes T};
\draw[onbo,line width=1.5pt,smooth]
  plot coordinates {(0.5,4) (2,3.6) (4,3) (6,2.2) (7.5,1.3) (9,0.5)};
\draw[onbmuted,dashed] (0.5,0) -- (0.5,4.1);
\node[below,font=\scriptsize] at (0.5,-0.1) {contamination};
\draw[onbmuted,dashed] (7.5,0) -- (7.5,1.3);
\node[below,font=\scriptsize] at (7.5,-0.1) {SIDA déclaré};
\node[above,onbo,font=\scriptsize] at (5,4.3) {phase asymptomatique (plusieurs années)};
\end{tikzpicture}
\fig{Figure — Évolution du nombre de lymphocytes T au cours de l'infection par le VIH.}
\end{illus}

\begin{remarque}
Grâce aux traitements antirétroviraux, une personne séropositive suivie médicalement
peut aujourd'hui maintenir un taux de lymphocytes T stable pendant très longtemps, mener
une vie normale, et réduire considérablement le risque de transmettre le virus.
\end{remarque}

\section{Les modes de transmission}

\begin{propriete}[Trois voies de transmission réelles]
\begin{itemize}
\item \textbf{voie sexuelle} : rapports sexuels non protégés ;
\item \textbf{voie sanguine} : partage de matériel d'injection, transfusion de sang non
contrôlé ;
\item \textbf{voie mère-enfant} : pendant la grossesse, l'accouchement ou
l'allaitement.
\end{itemize}
\end{propriete}

\begin{propriete}[Ce qui ne transmet PAS le VIH]
Le VIH \textbf{ne se transmet pas} par les gestes de la vie quotidienne : une poignée de
main, une embrassade, le partage d'un repas ou de couverts, une toux, un éternuement, une
piqûre de moustique, ou le fait de côtoyer une personne séropositive au travail ou à
l'école.
\end{propriete}

\begin{remarque}
Ces idées fausses sur la transmission sont à l'origine de nombreuses discriminations
envers les personnes séropositives, qui peuvent pourtant vivre, travailler et étudier
normalement aux côtés des autres, sans aucun risque de transmission par ces contacts
quotidiens.
\end{remarque}

\section{La prévention}

\begin{propriete}[Moyens de prévention]
\begin{itemize}
\item l'utilisation du \textbf{préservatif} lors des rapports sexuels ;
\item la \textbf{non-réutilisation} de matériel d'injection ;
\item le \textbf{dépistage régulier}, seul moyen de connaître son statut sérologique ;
\item le \textbf{traitement antirétroviral} précoce, qui protège la santé de la personne
infectée et réduit fortement le risque de transmission.
\end{itemize}
\end{propriete}

\begin{exemple}
Un couple qui souhaite connaître son statut sérologique avant une relation peut réaliser
un test de dépistage, simple, rapide et souvent gratuit, dans un centre de santé.
\end{exemple}

% ═══════════════════════════════════════════════════════════════
\section{L'essentiel à retenir}

\begin{synthese}
\textbf{VIH.} Virus qui détruit progressivement les lymphocytes T.\\[4pt]
\textbf{Trois phases.} Contamination/séroconversion $\to$ phase asymptomatique (souvent
plusieurs années) $\to$ SIDA déclaré (maladies opportunistes).\\[4pt]
\textbf{Transmission réelle.} Voie sexuelle, voie sanguine, voie mère-enfant.\\[4pt]
\textbf{Pas de transmission} par les gestes de la vie quotidienne (poignée de main, repas
partagé, moustique\dots).\\[4pt]
\textbf{Prévention.} Préservatif, matériel d'injection à usage unique, dépistage,
traitement antirétroviral.\\[4pt]
\textbf{Piège fréquent.} Une personne séropositive en phase asymptomatique se sent en
bonne santé mais peut transmettre le virus : l'absence de symptômes ne signifie pas
l'absence de risque.
\end{synthese}

% ═══════════════════════════════════════════════════════════════
\section{Méthodes \& savoir-faire}

\methode{Situer une phase de l'infection à partir d'une description}
Relier les symptômes (ou leur absence) et le taux de lymphocytes T décrits à l'une des
trois phases : contamination, phase asymptomatique, SIDA déclaré.
\begin{exemple}
Une personne séropositive sans symptôme depuis plusieurs années se trouve en phase
asymptomatique.
\end{exemple}

\methode{Distinguer un mode de transmission réel d'une idée fausse}
Vérifier si la situation correspond à l'une des trois voies réelles (sexuelle, sanguine,
mère-enfant) ou à un contact quotidien sans risque.
\begin{exemple}
Partager un repas avec une personne séropositive ne présente aucun risque de
transmission.
\end{exemple}

\methode{Proposer un message de prévention adapté}
Associer chaque moyen de prévention à la voie de transmission qu'il permet de réduire.
\begin{exemple}
Le préservatif réduit le risque lié à la voie sexuelle ; le matériel d'injection à usage
unique réduit le risque lié à la voie sanguine.
\end{exemple}

% ═══════════════════════════════════════════════════════════════
\section{Exercices d'application}
\begin{multicols}{2}

\rubrique{Le VIH et ses effets}
\exo{}\diff{1} Que signifie l'acronyme VIH ?

\exo{}\diff{2} Quelles cellules le VIH détruit-il progressivement ?

\exo{}\diff{2} Quelle en est la conséquence sur le système immunitaire ?

\rubrique{Évolution de l'infection}
\exo{}\diff{2} Citer, dans l'ordre, les trois phases de l'infection par le VIH.

\exo{}\diff{2} Qu'est-ce que la séroconversion ?

\exo{}\diff{2} Une personne séropositive en phase asymptomatique présente-t-elle des
symptômes ?

\rubrique{Transmission}
\exo{}\diff{2} Citer les trois voies réelles de transmission du VIH.

\exo{}\diff{2} Le VIH se transmet-il par une piqûre de moustique ? Justifier.

\rubrique{Prévention}
\exo{}\diff{1} Citer deux moyens de prévention contre le VIH.

\exo{}\diff{2} Pourquoi le dépistage est-il important ?

% ═══════════════════════════════════════════════════════════════
\end{multicols}

\section{Exercices d'entraînement}
\begin{multicols}{2}

\rubrique{Le VIH et ses effets}
\exo{}\diff{3} Expliquer pourquoi le VIH rend l'organisme vulnérable à des maladies
opportunistes qui ne sont normalement pas dangereuses pour une personne en bonne santé.

\rubrique{Évolution de l'infection}
\exo{}\diff{3} Pourquoi une personne séropositive en phase asymptomatique peut-elle
malgré tout transmettre le virus ?

\exo{}\diff{3} Expliquer, à l'aide d'un graphique donnant le nombre de lymphocytes T en
fonction du temps, comment reconnaître le passage au stade SIDA déclaré.

\exo{}\diff{2} Pourquoi un traitement antirétroviral précoce améliore-t-il le pronostic
d'une personne séropositive ?

\rubrique{Transmission}
\exo{}\diff{3} Un élève affirme qu'on peut attraper le VIH en serrant la main d'une
personne séropositive. Rédiger une réponse courte expliquant pourquoi c'est faux.

\exo{}\diff{2} Pourquoi le partage de matériel d'injection est-il un mode de transmission
du VIH ?

\exo{}\diff{3} Expliquer pourquoi la transmission mère-enfant peut être fortement réduite
par un suivi médical adapté pendant la grossesse.

\rubrique{Prévention et sensibilisation}
\exo{}\diff{3} Proposer un court message de sensibilisation destiné à rassurer une classe
sur les risques réels (et les idées fausses) liés à la présence d'un camarade
séropositif.

\exo{}\diff{2} Pourquoi la discrimination envers les personnes séropositives repose-t-elle
souvent sur une méconnaissance des modes de transmission ?

\exo{}\diff{3} Pourquoi le dépistage régulier est-il recommandé même en l'absence de
symptômes ?

% ═══════════════════════════════════════════════════════════════
\end{multicols}

\section{Sujets type BEPC}

\sujetbac[Comprendre l'infection par le VIH]
\begin{enumerate}[label=\arabic*)]
\item Quelles cellules le VIH détruit-il ?
\item Cite les trois phases de l'évolution de l'infection.
\item Que se passe-t-il lorsque le nombre de lymphocytes T devient trop faible ?
\item Pourquoi un traitement antirétroviral est-il important dès le début de
l'infection ?
\end{enumerate}

\sujetbac[Transmission et prévention]
\begin{enumerate}[label=\arabic*)]
\item Cite les trois voies réelles de transmission du VIH.
\item Cite trois situations de la vie quotidienne qui ne transmettent pas le VIH.
\item Cite deux moyens de prévention contre le VIH.
\item Pourquoi le dépistage est-il une étape essentielle de la prévention ?
\end{enumerate}

\sujetbac[Sensibilisation contre les préjugés]
\begin{enumerate}[label=\arabic*)]
\item Un camarade affirme qu'il est dangereux de partager la cantine avec un élève
séropositif. Rédige une réponse courte expliquant pourquoi cette affirmation est fausse,
en t'appuyant sur les modes de transmission réels du VIH.
\item Propose un conseil pratique pour vivre normalement aux côtés d'une personne
séropositive.
\item Rédige une phrase d'affiche de sensibilisation encourageant le dépistage.
\item Pourquoi est-il important de lutter contre les préjugés liés au VIH/SIDA ?
\end{enumerate}

% ═══════════════════════════════════════════════════════════════
\section{Corrigés}
\begin{multicols}{2}

\corrige{de l'exercice 1}
Virus de l'Immunodéficience Humaine.

\corrige{de l'exercice 2}
Les lymphocytes T.

\corrige{de l'exercice 3}
Il l'affaiblit progressivement, rendant l'organisme vulnérable aux infections.

\corrige{de l'exercice 4}
Contamination/séroconversion, phase asymptomatique, SIDA déclaré.

\corrige{de l'exercice 5}
L'apparition d'anticorps anti-VIH détectables dans le sang, quelques semaines après la
contamination.

\corrige{de l'exercice 6}
Non, elle se sent généralement en bonne santé, bien que le virus continue d'agir.

\corrige{de l'exercice 7}
Voie sexuelle, voie sanguine, voie mère-enfant.

\corrige{de l'exercice 8}
Non, le VIH ne se transmet pas par les piqûres d'insectes.

\corrige{de l'exercice 9}
Par exemple : l'utilisation du préservatif et le dépistage régulier.

\corrige{de l'exercice 10}
Car c'est le seul moyen de connaître son statut sérologique et, le cas échéant, de
débuter rapidement un traitement.

\corrige{de l'exercice 11}
Car le VIH détruit les lymphocytes T, cellules clés de l'immunité acquise : sans elles,
l'organisme ne peut plus se défendre efficacement, même contre des microbes habituellement
inoffensifs.

\corrige{de l'exercice 12}
Car le virus est présent dans son organisme et peut se transmettre par les voies
sexuelle, sanguine ou mère-enfant, même en l'absence de symptômes visibles.

\corrige{de l'exercice 13}
Le nombre de lymphocytes T devient très faible sur le graphique ; c'est ce seuil bas qui
marque le passage au stade SIDA déclaré, où apparaissent les maladies opportunistes.

\corrige{de l'exercice 14}
Car il permet de limiter la destruction des lymphocytes T et de maintenir un système
immunitaire fonctionnel plus longtemps.

\corrige{de l'exercice 15}
Le VIH ne se transmet pas par un contact comme une poignée de main : il nécessite un
contact avec du sang, des sécrétions sexuelles, ou une transmission mère-enfant.

\corrige{de l'exercice 16}
Car une aiguille contaminée peut mettre en contact direct le sang d'une personne
infectée avec celui d'une autre personne.

\corrige{de l'exercice 17}
Car un suivi médical permet un traitement antirétroviral adapté et des précautions lors
de l'accouchement, réduisant fortement le risque de transmission au bébé.

\corrige{de l'exercice 18}
Par exemple : « Le VIH ne se transmet pas par un contact quotidien : on peut étudier,
jouer et partager un repas avec un camarade séropositif sans aucun risque. »

\corrige{de l'exercice 19}
Car par peur ou ignorance des voies réelles de transmission, certaines personnes
associent à tort tout contact avec une personne séropositive à un risque.

\corrige{de l'exercice 20}
Car l'infection peut être présente sans symptôme pendant plusieurs années : seul un test
permet de la détecter tôt et de commencer un traitement adapté.

\corrige{du sujet 1}
1) Les lymphocytes T.\\
2) Contamination/séroconversion, phase asymptomatique, SIDA déclaré.\\
3) Le système immunitaire s'effondre et des maladies opportunistes peuvent se
développer.\\
4) Car il limite la destruction des lymphocytes T et préserve le système immunitaire
plus longtemps.

\corrige{du sujet 2}
1) Voie sexuelle, voie sanguine, voie mère-enfant.\\
2) Par exemple : une poignée de main, un repas partagé, une piqûre de moustique.\\
3) Par exemple : le préservatif et le dépistage régulier.\\
4) Car il permet de connaître son statut sérologique et de débuter un traitement
précocement si nécessaire.

\corrige{du sujet 3}
1) Le VIH ne se transmet pas par le partage d'un repas : cela nécessite un contact avec
le sang ou les sécrétions sexuelles, ou une transmission mère-enfant, ce qui n'est pas le
cas à la cantine.\\
2) Par exemple : continuer les activités habituelles (jeux, repas, travail scolaire)
normalement, sans crainte injustifiée.\\
3) Par exemple : « Se faire dépister, c'est se protéger et protéger les autres. »\\
4) Car ces préjugés isolent inutilement les personnes séropositives et peuvent dissuader
certaines personnes de se faire dépister par peur du jugement, ce qui nuit à la santé
publique.

\end{multicols}
\exercicesBanner
